\begin{landscape}
	\begin{table}[htbp]
		\centering
		\caption{Liste des 48 ressources classées en fonction de la phase de production.}
		\label{tab:ressources_production}
		
		\vspace{0.3cm}
		
		{\tiny
			\renewcommand{\arraystretch}{1.15}
			\setlength{\tabcolsep}{5pt}
			\setlength{\extrarowheight}{3pt}
			
			\begin{tabular}{|p{3.4cm}|p{6.8cm}|p{3.8cm}|p{3.2cm}|p{3.2cm}|}
				\hline
				\textbf{Pré-production (1)} \newline \textit{(production, préparation)} & 
				\textbf{Pré-production (2)} \newline \textit{(éducation, information)} & 
				\textbf{En-production} & 
				\textbf{Concerts} & 
				\textbf{Post-production} \\ 
				\hline
				
				• Déclaration d’intention \newline
				• Plan structurel du conducteur \newline
				• Concept structurel de la production \newline
				• Plan de répétition \newline
				• Document préparatoire divers \newline
				• Contenu dramaturgique \newline
				• Nomenclature & 
				
				• Conducteur original (source prim.) \newline
				• Partition originale (source prim.) \newline
				• Édition originale (source prim.) \newline
				• Livret original (source prim.) \newline
				• Nouvelle édition basée sur sources prim. (source sec.) \newline
				• Nouvelle partition basée sur nouvelle éd. (source sec.) \newline
				• Pitch original (source prim.) \newline
				• Tempérament original (source prim.) \newline
				• Illustration d’instrument (source prim./iconog.) \newline
				• Image d’instrument (photo/vidéo) (source sec./iconog.) \newline
				• Illustration d’ensemble, continuo, chef, position/posture (source prim.) \newline
				• Image d’ensemble, continuo, chef, position/posture (photo/vidéo) (source sec.) \newline
				• Extrait de traité instrumental (source prim.) \newline
				• Extrait de traité sur styles (source prim.) \newline
				• Extrait de livre musique hist. formée (source sec.) \newline
				• Image d’instrument de musée (source prim.) \newline
				• Style orig. d’interface utilisateur (archets, embouchures...) (source prim.) \newline
				• Instrument moderne ou interface conçue musique hist. formée (source sec.) \newline
				• Enregistrement (source prim.) & 
				
				• Comm. sur répertoire (rareté, pépite, fake news...) \newline
				• Plan de salle \newline
				• Ajustement (plan répétition, coupes, réorg., durée...) \newline
				• Compromis (instruments, pitchs, styles) \newline
				• Annotations de conducteur \newline
				• Annotations de partition \newline
				• Enregistrement de test \newline
				• Notes de programmes \newline
				• Photo répétitions (pupitres, interactions...) \newline
				• Interviews musiciens (audio, vidéo) \newline
				• Unpublished artist statement \newline
				• Vidéo matériel éducatif & 
				
				• Enregistrement final \newline
				• Photo concert / enregistrement public \newline
				• Matériel sortant (synthèse, kit presse...) \newline
				• Matériel éducatif de communication & 
				
				• Article de presse \newline
				• Partition avec annotations finales \newline
				• Bilan de production \newline
				• Évaluation plan de production \newline
				• Critique \newline
				• Conducteur avec annotations finales \\ 
				\hline
			\end{tabular}
		}
	\end{table}
\end{landscape}